% ============================================================
% SECTION 4 — The Flood Substrate
% ============================================================
\section{The Flood Substrate}
\label{sec:substrate}

\begin{figure*}[t]
  \centering
  \includegraphics[width=\textwidth]{figures/nfip_scenario_comparison.pdf}
  \caption{Cumulative NFIP claim footprint across the five CONUS benchmark
    scenarios, for scale and orientation. These are cumulative all-hazard
    claims (through 2024) within each scenario's geographic footprint
    ---\emph{not} the event-matched per-(\textit{zcta\_id},\textit{event})
    target the benchmark models (\S\ref{sec:problem-substrate}). Oahu is
    excluded here: it is an out-of-distribution probe
    (Appendix~\ref{app:oahu}), not a benchmark scenario.}
  \label{fig:nfip-scenarios}
\end{figure*}

The benchmark operates on a single multi-event flood substrate at the
(\textit{zcta\_id}, \textit{event}) grain, spanning five U.S.\ metropolitan
regions selected for complementary hazard profiles and data availability
(Table~\ref{tab:scenarios}). Houston provides the densest claim record
(396 ZCTA--event pairs) with both tropical and inland flooding; Southwest
Florida (606 pairs) covers hurricane surge and coastal inundation; New York
City (422 pairs) adds nor'easter and tidal flood events; Riverside--Coachella
Valley (172 pairs) contributes arid flash-flood events with thin insurance
penetration; and New Orleans (264 pairs) spans the pre- and post-levee
era, offering a natural experiment in infrastructure-mediated flood risk.

\begin{table}[t]
\centering
\caption{Benchmark scenarios. $n$ is ZCTA--event rows. Targets:
  \texttt{nfip} = NFIP event claims (regression); \texttt{311} = citizen
  reports (classification); \texttt{hwm} = USGS high-water marks
  (classification). A cell is modelable when the target has enough
  non-missing rows for spatial-blocked CV.}
\label{tab:scenarios}
\small
\begin{tabular}{@{}llcl@{}}
\toprule
\textbf{Scenario} & \textbf{$n$} & \textbf{Targets} & \textbf{Hazard profile} \\
\midrule
Houston         & 396 & nfip, 311, hwm & Tropical + inland \\
SW Florida      & 606 & nfip           & Hurricane surge \\
NYC             & 422 & nfip, 311, hwm & Nor'easter + tidal \\
Riverside       & 172 & nfip           & Flash flood (arid) \\
New Orleans     & 264 & nfip, hwm      & Pre/post-levee \\
\bottomrule
\end{tabular}
\end{table}

\subsection{Data layers}

Features are organized by representation level
(\S\ref{sec:method-ladder}), where each level is a strict superset of
the previous:

\paragraph{R0 (static tabular, 33--36 features).}
American Community Survey demographics~\cite{kwon2023acs}, CDC Social
Vulnerability Index, FEMA flood zone classifications, terrain (TWI,
slope, elevation from USGS 3DEP), HIFLD critical infrastructure counts,
and temporally-gated NFIP claim history (computed strictly before each
event's onset to prevent temporal leakage).

\paragraph{R1 (+hydrology and spatial structure, 58--63 features).}
R0 plus NHDPlus drainage network metrics, USACE levee proximity, estimated
sewershed capacity, and eight spatial-lag features derived from a
queen-contiguity adjacency matrix $\mathbf{W}$. The target spatial lag
(\textit{wlag\_nfip\_claims}) is recomputed per fold from training units
only to prevent test-fold leakage (\S\ref{app:relational-ablation}).

\paragraph{R2 (+temporal dynamics, 67--72 features).}
R1 plus NOAA MRMS Stage~IV quantitative precipitation estimates, HRRR
forecast fields, tide-gauge readings, and HURDAT2 storm-track proximity.
These sources have operational start dates (MRMS: October 2012; HRRR:
September 2014), so pre-dating events carry a
\texttt{not\_applicable\_temporal} field status rather than an imputed
value.

\subsection{Provenance and field status}

Each feature carries a provenance record (source agency, vintage,
spatial resolution) and one of four field-status values:
\texttt{available} (value present), \texttt{not\_applicable\_temporal}
(event predates the source's operational start),
\texttt{not\_applicable\_type} (source does not apply to this hazard
type), or \texttt{source\_empty} (source covers the event but returned
zero records). This vocabulary is the mechanism by which the benchmark
surfaces evidence gaps rather than silently imputing through them---the
allocation verdict, for instance, is necessarily partial because
temporal features carry \texttt{not\_applicable\_temporal} for pre-2012
events (Appendix~\ref{app:protocol-substrate},
Table~\ref{app:tab:availability}).

\subsection{Known limitations}

\paragraph{ZIP-as-ZCTA (C-001).} NFIP claims are reported by ZIP code;
ZCTAs are Census-defined approximations. Within the five metropolitan
footprints, fewer than 3\% of ZIPs map to a different ZCTA. We treat
ZIP as ZCTA and document the approximation, following standard
practice~\cite{krieger2002zip, chen2025crosswalk}. The mismatch rate is
higher in rural areas outside our study regions.

\paragraph{Riverside sample size.} At $n{=}172$ ZCTA--event pairs,
Riverside--Coachella Valley is the thinnest scenario and produces the
widest confidence intervals. Results from this scenario should be
interpreted with that constraint.
